\documentclass[10pt]{article}
\usepackage[margin=0.85in]{geometry}
\usepackage{amsmath,amssymb,amsthm,booktabs,hyperref,microtype}
\setlength{\emergencystretch}{1em}
\newcommand{\E}{\mathbb E}
\newcommand{\Prb}{\mathbb P}
\newcommand{\Var}{\operatorname{Var}}
\title{Review of the New Constant Proof: Analytic Checks and Audit Supplement}
\author{}\date{October 22, 2026}
\begin{document}\maketitle

\noindent\textbf{Verdict.} The sufficient guarantees and the necessary integer
constant conclusion in \path{results_parameter_search/proof.tex} survive
this review, under that manuscript's explicit model and quantifiers.
No mathematical counterexample or invalid implication was found. This is
an analytical review plus exact/ball arithmetic verification, not a proof
assistant formalization. Several compressed steps are expanded below.
The earlier independent numerical auditor did not separately check domain
completeness, counting-branch bounds, or final aggregation; the supplementary
auditor now checks those. This was an audit-coverage limitation, not a
discovered false inequality. Frozen proof, certificates, and experiments
have not been edited.

\section{What is and is not established}
The statistic is the original multi-scale Chebyshev statistic, with
$m=Cn/(\epsilon(1+\log n))$, independent Poisson counts, no sample cap,
no polynomial truncation, and exact real arithmetic. Here $n$ ranges over
all positive integers and $\epsilon$ over $(0,1)$. The distribution may
have countable support. On the binary promises $S\le n$ and
$\Delta_n\ge\epsilon$, the two outputs have simultaneous correctness
with the claimed probability. In the gap, only the support-bound guarantee
is asserted.

\begin{center}
\begin{tabular}{rrr}
\toprule $C$ & $L_0$ & strict single-run failure bound\\\midrule
70&14&$0.0009$\\
70&15&$0.000081$\\
70&16&$0.000057$\\
70&18&$0.000040$\\\bottomrule
\end{tabular}
\end{center}
The necessary result is about this unchanged construction with fixed
finite $L_0>1$. It is not a lower bound for every support tester, for a
modified score, or for a restricted finite experimental domain.
It establishes $C\ge C_{\rm crit}$ as a necessary condition, where
\[
69.00627138<C_{\rm crit}<69.00627139.
\]
The sufficient $(70,14)$ pair therefore proves the smallest admissible
\emph{integer} $C$ is 70. At $C=C_{\rm crit}$ the argument makes no
claim. More precisely, the infimum of admissible real constants lies in
$[C_{\rm crit},144(\log2)^2]$; attainment of that infimum is not proved.
The latter endpoint is approximately $69.185234004221$.
Neither the minimal $L_0$ nor the fastest practical implementation is
determined. Fourteen is a sufficient branch threshold.

\section{Concentration: range, centering, and dependence}
Let $I_j=\sum_{\ell=0}^{k_j}|c_{j,\ell}|\ell!/m^\ell$, with constant
term of absolute value one. Alternating coefficient signs give both
$\|f_j\|_\infty\le I_j$ and range length at most $I_j$.
The scale-zero coefficient envelope $F_{\rm coef}$ also bounds absolute
value, hence its range is at most $2F_{\rm coef}$. Consequently
\[
\operatorname{range}(f_0)^2\le
4\min(I_0^2/4,F_{\rm coef}^2)=4V_0.
\]
This point matters: a variance upper bound alone would not justify the
Hoeffding term. Here the squared-range upper bound is actually available.
For the other scales $V_j=I_j^2/4$ has the same property.

Put $\delta=\delta_j$, $\bar\rho=\delta^2/2$, $d=2-\delta/2$.
If $S\le n$, the upper-error threshold has gap
$x=\delta n+(\delta-\bar\rho)S$ above $(1+\bar\rho)S$.
For $u=S/n\in(0,1]$,
\[
\frac{S}{x^2}=\frac{1}{n\delta^2}
\frac{u}{(1+(1-\delta/2)u)^2}
\le\frac{1}{n\delta^2d^2},\qquad
\frac1x\le\frac1{\delta n}.
\]
The derivative of the fraction is nonnegative on $[0,1]$.
This validates the reduction used for both Cantelli and Hoeffding.

For the new endpoint Bernstein bound, set $G=I_j-1>0$ and
$A=1+\bar\rho+\max(0,G-1)$. The variable
$Y=f_j(N)-(1+\bar\rho)$ has $\E Y\le0$ and lies in $[-A,G]$.
For $t\ge0$, the chord inequality is
\[
e^{ty}\le\frac{G-y}{A+G}e^{-tA}
             +\frac{A+y}{A+G}e^{tG}.
\]
The right side's expectation is increasing in $\E Y$, so replacing it
by zero only increases the bound. The resulting two-point distribution
has mean zero and variance $AG$. Independence across labels thus gives
\[
\Prb\{\textstyle\sum_{i=1}^S Y_i\ge x\}
\le \exp\left\{-\frac{x^2}{2SAG+2Gx/3}\right\}
\le \exp\left\{-\left[
\frac{2AG}{n\delta^2d^2}+\frac{2G}{3\delta n}\right]^{-1}\right\}.
\]
It is legitimate to take the minimum with Cantelli and Hoeffding since
each bounds the same event. It is also legitimate to sum over scales
by a union bound. No independence between scales is assumed or needed.

For the lower error, $f_j(0)=0$ implies
\[
\sum_i\Var f_j(N_i)\le F_j^2\E D.
\]
If $\E D\le\alpha n$, $\alpha=33/32$, the centered lower-tail summands
are at most $2F_j$. Bernstein at the expectation margin
$\kappa_j\delta_j n$ gives exactly
\[
\exp\left\{-\frac{\kappa_j^2}
 {2\alpha w_j+(4/3)\kappa_j\sqrt{w_j/n}}\right\},
\quad w_j=\frac{F_j^2}{n\delta_j^2}.
\]
If $\E D>\alpha n$, the independent Bernoulli occupancy indicators give
$\Prb(D\le n)\le\exp[-n(\alpha-1-\log\alpha)]\le e^{-n/2112}$.
The inherited heavy-mass, small-mass, and intermediate-mass expectation
arguments are exhaustive cases for a fixed distribution. Taking a maximum
of their bounds, rather than adding them, is correct.

\section{Countable support and simultaneous outputs}
Here is an explicit justification of the manuscript's passage to infinite
sums. For a fixed scale and input, $F_j<\infty$ and $m<\infty$.
As $\sum_i\Prb(N_i>0)=\E D\le m$,
\[
\sum_i\E|f_j(N_i)|\le F_jm<\infty.
\]
Tonelli implies absolute almost-sure convergence. In fact the total
sample count is finite almost surely, so only finitely many scores are
nonzero. The sum of expectations is absolutely convergent as well.
For a tail set of labels of total probability $\tau$, let $D_{\rm tail}$
be its occupancy count. Independence gives
\[
\E D_{\rm tail}^2\le m\tau+(m\tau)^2,\qquad
\E\left|\sum_{\rm tail}f_j(N_i)\right|^2
 \le F_j^2[m\tau+(m\tau)^2]\longrightarrow0.
\]
Thus partial sums also converge in $L^2$. Centered finite partial sums
have the Bernstein moment-generating bound with a common finite variance
proxy. Apply Fatou to their nonnegative exponentials and then exponential
Markov to obtain the same bound for the limit. This avoids any unsupported
claim that an infinite sum of range lengths is finite: Hoeffding is only
used in the finite-support case $S\le n$.

Write $t=\min(e_\epsilon,n)$. If $t>1$, then
$\Delta_{t-1}>\epsilon$. Hence the lower-error bound applies at $r=t-1$.
For $S\le n$, the event
$\widetilde S\le S$ and $\widetilde S>t-1$ ensures both correct outputs;
when $t=1$ the lower constraint is automatic. This costs at most $U+R_f$.
Under the alternative, $\widetilde S>n$ implies rejection and
$\widehat S=n$, so the single bound $R$ ensures both outputs. In the gap
the upper support bound is automatic since $S>n$ and only the lower
constraint is required. Thus the common bound
$\max(U+R_f,R)$ is correct, including equality conventions.

\section{Why the asymptotic obstruction really excludes 69}
Fix $C\ge1$ and write $L=1+\log n$,
$a=\sqrt{512/C}\,3\log2$, $k=\lceil aL\rceil$,
$A=C/(8L)$, and $B=64L$. These parameters, and the dimensionless
factorial weights $u_\ell$, do not depend on $\epsilon$.
The exact integral identity is
\[
I_L=\int_0^\infty e^{-s}
\frac{T_k(z_L+2s/(B-A))}{T_k(z_L)}\,ds,
\qquad z_L=(B+A)/(B-A).
\]
We have $k\operatorname{arcosh}z_L\to6\log2$.
For $s=yL$, uniformly on compact subsets of $y>0$,
the logarithm of the integrand divided by $L$ converges to
\[
\phi(y)=-y+a\operatorname{arcosh}(1+y/32).
\]
To justify uniformity explicitly, the arguments of arcosh converge
uniformly to a compact subset of $(1,\infty)$, $k/L\to a$, and
$|\log\cosh z-z|\le\log2$ for $z\ge0$. The denominator's logarithm
is $O(1)$. The positive maximizer and maximum are
\[
y_* = \sqrt{1024+a^2}-32,\qquad
\gamma(C)=a\operatorname{arsinh}(a/32)-\sqrt{1024+a^2}+32.
\]
For any $\eta>0$, choose a compact interval $K$ about $y_*$ where
$\phi\ge\gamma-\eta/3$. Uniform convergence gives integrand at least
$\exp[(\gamma-2\eta/3)L]$ there for large $L$. The substitution contributes
the prefactor $L|K|$, which is at least $e^{-\eta L/3}$ eventually.
Therefore $I_L\ge e^{(\gamma-\eta)L}$. Only this lower bound is needed.

Choose a maximizing $\ell_L\in\{1,\ldots,k\}$ and let
\[
p_n=(1/2,1/[2(n-1)],\ldots,1/[2(n-1)]),\quad
\epsilon_n=\frac{Cn}{2(n-1)\ell_LL}.
\]
This is a null distribution with exactly $n$ labels. The parameter choice
is not circular because the weights determining $\ell_L$ are independent
of $\epsilon$. Moreover $0<\epsilon_n\le C/L<1$ eventually, and the
prescribed sample mean is $m=2(n-1)\ell_L$. The rare-label counts are
independent Poisson variables of integer mean $\ell_L$.

Let $X_L=f_0(\operatorname{Poi}(\ell_L))$, with mean $\mu_L$ and
variance $\sigma_L^2$. The expectation lemma applies eventually because
$a(C)<64$ for $C\ge1$, $\sqrt{B(B-A)}/L\to64$, and $b\to0$.
It gives $0\le\mu_L\le33/32$; also $|X_L|\le I_L$.
The integer mode of a Poisson variable satisfies
$\Prb(\operatorname{Poi}(\ell)=\ell)\ge3/(20\sqrt\ell)$:
Chebyshev puts at least $3/4$ of the mass in at most $5\sqrt\ell$
integers, and the mode dominates their individual masses.
Since $u_{\ell_L}\ge(I_L-1)/k$, eventually
$|f_0(\ell_L)|\ge I_L/(4k)$. Thus
\[
\sigma_L^2\ge\frac{3I_L^2}{320k^{5/2}}-(33/32)^2
\ge c\frac{I_L^2}{k^{5/2}}.
\]
Take $s_n=\sqrt{n-1}\,\sigma_L$. The centered rare scores satisfy
\[
\max_i\frac{|X_{n,i}-\mu_L|}{s_n}
 \le O\left(\frac{k^{5/4}}{\sqrt n}\right)\to0,
\qquad \sum_{i=2}^n\Var\frac{X_{n,i}-\mu_L}{s_n}=1.
\]
The Lindeberg sum is therefore exactly zero eventually for every fixed
positive cutoff. This explicitly verifies the triangular-array CLT
hypothesis; an ordinary fixed-distribution iid CLT would be insufficient.
The reference theorem was rechecked in
\href{https://people.eecs.berkeley.edu/~jordan/courses/210B-spring07/lectures/stat210b_lecture_2.pdf}
{Jordan--Kleiner, Stat210B Lecture 2, Theorem 9}.

If $\gamma(C)>1/2$, choose $0<\eta<\gamma(C)-1/2$. Then
\[
\frac n{s_n}\le O\left(\frac{\sqrt n\,k^{5/4}}{I_L}\right)
\le O\left(L^{5/4}e^{-(\gamma-\eta-1/2)L}\right)\to0.
\]
The head count has mean $\lambda_n=(n-1)\ell_L\ge n-1$.
For large $n$, Chebyshev bounds
$\Prb(N_1\le k)\le\lambda_n/(\lambda_n-k)^2\to0$,
so its score is one with probability tending to one.
The standardized threshold
$(3n/2-1-(n-1)\mu_L)/s_n$ tends to zero. Combining the CLT and the
vanishing probability of an exceptional head score proves
\[
\Prb(Z_0>3n/2)\to1/2.
\]
Because $\delta_0=1/4$, this event implies $\widetilde S>n$.
For the full maximum statistic the justified statement is
$\liminf\Prb(\text{false rejection})\ge1/2$, not necessarily equality.
Every fixed finite $L_0$ eventually activates the polynomial branch.

The derivative of $\gamma$ with respect to $a$ is
$\operatorname{arsinh}(a/32)>0$, while $a$ decreases with $C$.
The unique root of $\gamma=1/2$ has the bracket stated above;
$\gamma(69)\approx0.500045328>1/2$.
The strict excess is small, and this proof gives no useful finite-size
failure threshold for 69. Finite successful experiments do not contradict it.
For $0<C\le1$, the $n=1$ counting-branch alternative in the manuscript
has limiting error $(1+C)e^{-C}\ge2/e>1/2$.
For any fixed odd repetition count, the majority-error polynomial is
continuous, increasing, and takes value $1/2$ at input $1/2$.
Thus fixed odd majority repetition cannot remove either uniform obstruction.

\section{Arithmetic and domain verification}
The review checks the inherited expectation, integral/range, coefficient,
and half-line inequalities, including their side conditions and the
directions of monotonic substitutions. On $14\le L\le40$, the 260
closed slabs for each constant cover every possible real $L$ without gaps.
The saved scale vectors include every active $j$, and
$k_j=\lceil\alpha_jL\rceil\le\alpha_jL_{\rm right}+1$ includes every
degree transition. For $L\ge40$, the finite low-scale envelopes and the
geometric omitted-scale sum have the required derivative signs. The
equality $2t_0=1$ at $C=144(\log2)^2$ is treated symbolically.

In the counting branch, before the $(r+1)$st discovered label the unseen
probability mass is at least $\Delta_r\ge\epsilon$. Conditional exponential
waiting times can therefore be coupled with rate-$\epsilon$ waiting times.
This gives
$\Prb(D\le r)\le\Prb(\operatorname{Poi}(m\epsilon)\le r)
\le\Prb(\operatorname{Poi}(Cn/L)\le n)$.
The quarter-width slab Chernoff bounds have now been independently
recomputed. Their largest certified bounds are approximately
$2.580270\cdot10^{-12}$ for $C_*$ and $1.862823\cdot10^{-12}$ for 70.
These are safely below $10^{-10}$.

The supplementary program \texttt{parameter\_search/review\_proof.py}
also verifies exact rational equality between each saved final upper
endpoint and the maximum over its counting bound, selected slabs, and
half-line bound. It checks the critical root using the positive-series
identity
\[
\operatorname{arsinh}x=2\operatorname{artanh}z,
\quad z=\frac{x}{1+\sqrt{1+x^2}},\quad
0\le 2\operatorname{artanh}z-2\sum_{k=0}^{N-1}
 \frac{z^{2k+1}}{2k+1}
\le\frac{2z^{2N+1}}{(2N+1)(1-z^2)}.
\]
This differs from the original alternating asinh series. At 512-bit Arb
precision it confirms both root endpoints, exclusion of 69, and the
opposite sign at 70.

The supplementary audit passes 14,880 explicit checks; the separate
existing Arb formula audit passes 20,507 checks. These counts are code
obligations, not counts of independent mathematical theorems. A fresh
isolated rerun of the original exact certificate and root certificate is
recorded in \texttt{reproduction\_comparison.json}; the rerun does not
overwrite the frozen artifacts. The exact verifier checks 32,752
obligations using outward dyadic/rational arithmetic. The check scripts
cannot by themselves establish the CLT, infinite-sum limits, or the
analytical reduction; those arguments are reviewed explicitly above.

\section{Practical interpretation}
The stronger result changes the certified error analysis, not the estimator
or its per-run sample mean. At the same $(C,L_0)$, one run is the same
mathematical procedure as before. Its older loose error bound requiring
many repetitions has been replaced by a sharper bound.
The theorem covers an ideal Poissonized statistic. An implementation's
degree cap, resource limits, floating-point input domain, or abort policy
needs its own correctness statement; those restrictions are not silently
included in this theorem. Sample savings do not establish CPU-time savings.
The supported conclusion is integer optimality of $C$ within the stated
construction and uniform domain, together with certified sufficient $L_0$
values. Joint global optimality of both parameters remains unproved.
\end{document}
